\\documentclass[11pt, a4paper]{article} %tamaño mínimo de letra 11pto.

\usepackage{graphicx} 
%\usepackage[english]{babel} %Español 
%\usepackage[utf8]{inputenc} %Para poder poner tildes
\usepackage{vmargin} %Para modificar los márgenes
\setmargins{2.5cm}{1.5cm}{16.5cm}{23.42cm}{10pt}{1cm}{0pt}{2cm}
%margen izquierdo, superior, anchura del texto, altura del texto, altura de los encabezados, 
% entre el texto y los encabezados, altura del pie de página, espacio entre el texto y el pie de página

% introducido por mi 
\usepackage{amsmath,amsfonts,amssymb} % ecuaciones
\usepackage{siunitx}
\usepackage{fancyhdr} % headers y footers
\usepackage{esint} % integrales cerradas
\usepackage{bbold} % para meter matriz identidad
\usepackage{graphicx} % estos dos para una figura
\usepackage{caption}
\usepackage{subcaption}
\usepackage{pdfpages} % para incluir declaracion IA
\usepackage{float}
%\usepackage[numbers,sort&compress]{natbib} % for bibliography
\usepackage[inkscape=false]{svg} % add svg files
\usepackage{hyperref} % to add video
\usepackage{tikz} % to generate plots 
\usepackage{setspace} % to space out text lines
\usepackage[T1]{fontenc}
\usepackage{tikz}
%\usepackage{standalone}
\newcolumntype{C}{>{$}c<{$}}

\usepackage[style=numeric,sorting=none, backend=biber]{biblatex}  % Load biblatex with the 'verbose' style for full citations in footnotes
\addbibresource{bibliography.bib}

\setlength{\parskip}{\baselineskip}
\renewcommand{\baselinestretch}{1.2}

\newcommand{\xv}{\vec{x}}
\newcommand{\yv}{\vec{y}}
\newcommand{\xd}{\dot{x}}
\newcommand{\xdd}{\ddot{x}}
\newcommand{\yd}{\dot{y}}
\newcommand{\ydd}{\ddot{y}}
\newcommand{\bigO}{\mathcal{O}}
\newcommand{\Th}{\theta}

% \usetikzlibrary{decorations.markings,arrows.meta} % for big arrows
% \usetikzlibrary{calc}
% \tikzset{
%     >={Stealth[length=3pt]}, % Standard, clean arrowhead
%     % Custom style to place a single arrow at the halfway point of a path
%     cleanflow/.style={
%         decoration={
%             markings,
%             mark=at position 0.5 with {\arrow{>}}
%         },
%         postaction={decorate}
%     },
%     stable/.style={line width=1.5pt, draw=black},
%     thinarc/.style={line width=0.7pt, draw=black!90},
%     textlabel/.style={font=\small\selectfont, text=black}
% }

% \tikzset{
%   flow/.style={
%     decoration={
%       markings,
%       mark=at position #1 with {\arrow[scale=1.3]{stealth}}
%     },
%     postaction={decorate}
%   },
%   flow/.default=0.55
% }

% % Layer dimensions
% \newcommand{\layerW}{3.0}   % width
% \newcommand{\layerH}{0.45}  % height of magnetic layers
% \newcommand{\barrH}{0.30}   % height of tunnel barrier

% % Colors
% \definecolor{ferrocolor}{RGB}{173, 216, 230}   % light blue – ferromagnet
% \definecolor{barrcolor}{RGB}{255, 228, 196}    % bisque – insulating barrier

% \tikzset{
%   mtjlayer/.style={draw=black, thick, minimum width=\layerW cm, minimum height=\layerH cm, inner sep=0},
%   barrier/.style={draw=black, thick, minimum width=\layerW cm, minimum height=\barrH cm,
%                   inner sep=0, fill=barrcolor},
%   magarrow/.style={-{Stealth[length=6pt, width=4pt]}, ultra thick, black},
% }

% \newcommand{\drawMTJ}[4]{
%   % bottom ferromagnet (free layer)
%   \node[mtjlayer, fill=ferrocolor] at (#1, 0) {};
%   % tunnel barrier
%   \node[barrier] at (#1, \layerH*0.5 + \barrH*0.5) {};
%   % top ferromagnet (reference layer)
%   \node[mtjlayer, fill=ferrocolor] at (#1, \layerH + \barrH) {};

%   % Arrow in top (reference) layer
%   \draw[magarrow]
%     ({#1 - 0.55*#2}, {\layerH + \barrH})
%     -- ({#1 + 0.55*#2}, {\layerH + \barrH});

%   % Arrow in bottom (free) layer
%   \draw[magarrow]
%     ({#1 - 0.55*#3}, 0)
%     -- ({#1 + 0.55*#3}, 0);

%   % Layer labels
%   \node[right, font=\small] at ({#1 + \layerW*0.5}, {\layerH + \barrH}) {Fixed};
%   \node[right, font=\small] at ({#1 + \layerW*0.5}, {\layerH*0.5 + \barrH*0.5}) {Barrier};
%   \node[right, font=\small] at ({#1 + \layerW*0.5}, 0) {Free};

%   % Bit label below
%   \node[font=\large\bfseries] at (#1, -0.6) {#4};
% }

% \newcommand{\yinyanglobe}[3]{
%     \begin{scope}[rotate=#1]
%         \fill[#2] (0,0) 
%             arc[start angle=-90, end angle=90, radius=0.5cm] 
%             arc[start angle=90, end angle=-90, radius=1cm]
%             arc[start angle=90, end angle=270, radius=0.5cm] -- cycle;
            
%         \fill[#3] (0, 1) circle (0.25cm);
%     \end{scope}
% }

\begin{document}
Hello

\subsection{Theory}
In this section we consider the simplest model that captures the essence of non-hermitian behaviour: 2 linear actuators 
coupled non-reciprocally. We denote the two degrees of freedom as $x$ and $y$, which correspond to the displacement of the 
actuators from the equilibrium position. We choose the specfici non-reciprocal interaction to be odd elasticity. 
One can interpret the system as a 2D oscillator. 
The equations of motion are
    \begin{align}
        \xdd +k_0x +  \gamma \xd + k_cx^3 + k_ay = 0 \label{eq_1osc_x}\\ 
        \ydd +k_0y +  \gamma \yd + k_cy^3 - k_ax = 0 \label{eq_1osc_y}
    \end{align}
[should I add nonlinear cubic term to stabilize??]
where \( k_0 \) is the harmonic bound strength, \( \gamma \) is the damping coefficient, $k_c$ is a nonlinear saturation coefficient 
that stabilizes the dynamics, and \( k_a \) is the non-reciprocal coupling strength, also called the activity or gain. 
%The parameter \( \epsilon \ll 1 \) indicates that both damping and non-reciprocal coupling are weak perturbations to the harmonic oscillator.
There are different ways to analyse the system. One way is to plug in the plane wave ansatz 
\( x(t) = X e^{-iEt} \) and \( y(t) = Y e^{-iEt} \) in the linearized equations, and solve for the frequencies/energies $E$. 
Another method is to perform the multiple scale expansion, as explained in the previous section.
    The energies are 
    \begin{align}
        \text{clockwise:} \; E_1 = -\frac{i}{2}\left(\gamma + \frac{k_a}{\sqrt{k}}\right) \label{eq_1osc_e1}\\
        \text{counterclockwise:} \; E_2 = \frac{i}{2}\left(-\gamma + \frac{k_a}{\sqrt{k}}\right) \label{eq_1osc_e2}
    \end{align}
    and the eigenvectors are 
    \begin{align}
        \text{clockwise} \; v_1 = \begin{pmatrix}
        1 \\ -i
        \end{pmatrix} \label{eq_1osc_v1}\\
        \text{counterclockwise} \; v_2 = \begin{pmatrix}
        1 \\ i
        \end{pmatrix} \label{eq_1osc_v2}.
    \end{align}
    On the one hand, $E_1$ has 2 eigenvectors in which the oscillator rotates clockwise in the plane. 
    What distinguishes these eigenvectors is their initial phase, basically a global phase. As the time evolution is given by the 
    phase factor $e^{-iEt}$, these modes are exponentially suppressed as time evolves.
    On the other hand, $E_2$ has 2 eigenvectos that rotate counterclockwise, each with a different initial phase. 
    There are two regimes for the system described by equations \eqref{eq_1osc_x} and \eqref{eq_1osc_y}. 
    If $|k_a| < \gamma \sqrt{k}$, so for low values of activity, the energies have all a negative imaginary part. 
    This means that the oscillations of the actuators are exponentially damped, because friction dominates. The origin is a stable fixed point, 
    so the system becomes at rest and there is no motion. However, if the activity $|k_a|$ is increased, the system undergoes a supercritical 
    Hopf bifurcation. For $|k_a| > \gamma \sqrt{k}$, the origin is unstable and a stable limit cycle appears. The cubic 
    term $k_cx^3$ is necessary for the limit cycle to be stable. The sign of the activity $k_a$ determines the direction of 
    the limit cycle. If $k_a>0$, the limit cycle has counterclockwise direction, and if $k_a<0$, it has clockwise direction. 
    This feature can be seen in equations \eqref{eq_1osc_e1} and \eqref{eq_1osc_e2} by changing the sign of $k_a$. If the sign is changed, 
    the modes that can have the limit cycles, only for sufficiently large activity, are the ones corresponding to the clockwise chirality. 

    By plotting the trajectory of the system in state space, i.e. in the $x$-$y$ plane, we can see the effect of the non-reciprocal coupling.
    The system traces out limit cycles above the Hopf bifurcation. Limit cycles are isolated closed trajectories. 
    The fact that they are closed means that they are periodic solutions. They are usually defined in phase space, 
    and they are best visualized when the phase space is two-dimensional. However, our system 
    of equations \eqref{eq_1osc_x} and \eqref{eq_1osc_y} has actually four degrees of freedom: 
    the displacements $x$ and $y$, and their velocities $\xd$ and $\yd$. The phase space is thus four-dimensional.
    However, we can project the trajectories onto the $x$-$y$ plane, and we will see that they also trace out limit cycles there.
    Systems with negative chirality trace out limit cycles in clockwise direction, 
    while systems with positive chirality trace out limit cycles in counterclockwise direction. 
    This system's limit cycles are vey specific. From the eigenvector formulas \eqref{eq_1osc_v1} and \eqref{eq_1osc_v2}, 
    we can observe that the two linear actuators differ in value by a constant factor $\pm i = e^{\pm i\pi/2}$. 
    It means that the two actuators have a phase difference of $\pm \pi/2$, and they are said to be in quadrature. 
    If the actuators have the same amplitude, the limit cycles in this particular case are circles. If the actuators 
    have the same amplitude but other phase difference, they will trace out ellipses. 

    In conclusion, we use odd elasticity (or non-reciprocity in general) between two linear actuators to achieve stable limit cycles 
    that can oscillate either in clockwise or counterclockwise direction, and they can be used to replace the spin. We refer to the 
    mechanical analogue of spin as oscillator chirality or mechanical spin. 

% \begin{figure}[htbp]
%     \centering
%     % Suppress automatic sub-labels (a), (b), etc. directly below the panels
%     \captionsetup[subfigure]{labelformat=empty}
    
%     % =========================================================================
%     % LEFT COLUMN: 3 SUBFIGURES STACKED VERTICALLY
%     % =========================================================================
%     \begin{minipage}[b]{0.48\textwidth}
%         \centering
        
%         % Subfigure 1 (Top Left)
%         \begin{subfigure}[b]{\textwidth}
%              \begin{tikzpicture}[scale=0.4]
%     \useasboundingbox (-2.2, -2.5) rectangle (2.2, 1.5);
%         %\draw[red, dashed] (current bounding box.south west) rectangle (current bounding box.north east);
%   % Trajectory from bottom-left, spirals into fixed point
%   \draw[thick, flow=0.45]
%     (-2.0,-1.0) .. controls (-1.5,-0.8) and (-1.2, 0.4) ..
%     (-0.9, 0.55) .. controls (-0.5, 0.75) and ( 0.1, 0.5) ..
%     ( 0.25, 0.15) .. controls ( 0.35,-0.1) and ( 0.2,-0.3) ..
%     ( 0.0,-0.25) .. controls (-0.15,-0.1) and (-0.05, 0.0) ..
%     (-0.02, 0.0);

%   % Second trajectory from top-right, spirals into fixed point
%   \draw[thick, flow=0.45]
%     ( 2.0, 1.0) .. controls ( 1.5, 0.8) and ( 1.2,-0.4) ..
%     ( 0.9,-0.55) .. controls ( 0.5,-0.75) and (-0.1,-0.5) ..
%     (-0.25,-0.15) .. controls (-0.35, 0.1) and (-0.2, 0.3) ..
%     ( 0.0, 0.25) .. controls ( 0.15, 0.1) and ( 0.05, 0.0) ..
%     ( 0.02, 0.0);

%   % Stable fixed point (filled dot)
%   \filldraw (0,0) circle (2pt);

%   % Label
%   \node at (0,-2.4) {1};

%             \end{tikzpicture}
%             Hopf
% \begin{tikzpicture}[scale=0.4]
%         \useasboundingbox (-2.2, -2.5) rectangle (2.2, 1.5);
%         %\draw[red, dashed] (current bounding box.south west) rectangle (current bounding box.north east);
%   % Stable limit cycle (thicker, radius 1 cm)
%   \draw[line width=2pt, flow=0.25] (0,0) circle (1.0cm);

%   % Trajectory spiralling outward from near origin to limit cycle
%   \draw[thick, flow=0.5]
%     (0.15, 0.0)
%     .. controls ( 0.18, 0.30) and (-0.25,  0.55) ..
%     (-0.5,  0.35)
%     .. controls (-0.80, 0.15) and (-0.90, -0.35) ..
%     (-0.5, -0.80);   % ends near the circle r ≈ 1

%   % Trajectory spiralling inward from outside to limit cycle
%   \draw[thick, flow=0.45]
%     ( 1.80,-0.60)
%     .. controls ( 1.40,-1.60) and ( 0.40,-2.00) ..
%     (-0.60,-1.60)
%     .. controls (-1.40,-1.20) and (-1.30,-0.20) ..
%     (-1.00, 0.30);   % ends near the circle r ≈ 1

%   % Unstable fixed point (open dot)
%   \filldraw[fill=white, draw=black, thick] (0,0) circle (2pt);

%   % Label
%   \node at (0,-2.4) {2};

% \end{tikzpicture}
%             \caption{}
%             \label{fig:panel_a}
%         \end{subfigure}
        
%         %\vspace{0.5cm} % Vertical gap
        
%         % Subfigure 2 (Middle Left)
%         \begin{subfigure}[b]{\textwidth}
%             \includegraphics[width=0.7\textwidth]{Figures one unit/limit_cycle_example.pdf}
%             \caption{}
%             \label{fig:panel_b}
%         \end{subfigure}
        
%     \end{minipage}
%     \hfill
%     % =========================================================================
%     % RIGHT COLUMN: 2 SUBFIGURES STACKED VERTICALLY
%     % =========================================================================
%     \begin{minipage}[b]{0.48\textwidth}
%         \centering
        
%         % Subfigure 4 (Top Right)
%         \begin{subfigure}[b]{\textwidth}
%             \centering
%             \includegraphics[width=0.95\textwidth]{Figures one unit/global_scaling_profiles_stacked.pdf}
%             \caption{}
%             \label{fig:panel_d}
%         \end{subfigure}
        
%         \vspace{0.8cm} % Adjust gap to vertically balance against the left column
        
%         % Subfigure 5 (Bottom Right)
%         \begin{subfigure}[b]{\textwidth}
%             \centering
%             \includegraphics[width=0.95\textwidth]{Figures one unit/amplitude_scaling_Hopf.pdf}
%             \caption{}
%             \label{fig:panel_e}
%         \end{subfigure}
%     \end{minipage}

%     % =========================================================================
%     % MAIN UNIFIED CAPTION
%     % =========================================================================
%     \caption{Comprehensive structural and dynamical analysis of the classical altermagnetic analog lattice. Left column: (a) Structural topology of the active lattice units, (b) corresponding transient phase slip profile, and (c) additional left-side metric. Right column: Extracted global numerical scaling trends showcasing (d) limit cycle amplitude variations and (e) linear frequency scaling response under continuous activity sweeps.}
%     \label{fig:master_thesis_5_panel_compiled}
% \end{figure}



    \section{Experiments}

    The regime below the Hopf bifurcation, where the system goes to rest, is not interesting, so we focus on the regime 
    above the Hopf bifurcation where the system traces out limit cycles. There is only one parameter to control in the experiments, 
    and that is the activity $k_a$. The rest of coupling coefficients in equations \eqref{eq_1osc_x} and \eqref{eq_1osc_y} are 
    determined by the actuators' mechanical properties, which we don't change. We can decompose the analysis of the limit cycles into two: 
    the phase and the amplitude. 


\end{document}